%============================================================
\chapter{Khiêm Tốn Và Trí Tuệ (Humility and Wisdom)}
\begin{center}
  \textit{\color{truyenblue}The wisest man is he who knows he knows nothing.}\\
  \textit{(Người khôn nhất là kẻ biết rằng mình chưa biết gì.)}\\[4pt]
  \small\color{truyengold}--- Socrates ---
\end{center}
\ngancach

\section{Khổng Tử Và Cậu Bé Bảy Tuổi}
\begin{truyen}{Khổng Tử Hỏi Đường Trẻ Thơ}{Trung Hoa Cổ Đại}
\chuhoa{K}{}K hổng Tử cùng các học trò đang đi xe qua một làng, xe bị một đứa bé chắn ngang đường.\\
\textit{(Confucius and his disciples were traveling through a village when a young boy blocked their cart.)}

Đứa bé đang xây thành bằng đá cuội giữa đường và không chịu nhường.\\
\textit{(The boy was building a stone fortress in the middle of the road and refused to give way.)}

Khổng Tử xuống xe hỏi: ``Cháu không biết là phải nhường đường cho xe sao?''\\
\textit{(Confucius stepped down and asked: ``Don't you know you should give way to a cart?'')}

Cậu bé đáp: ``Từ xưa đến nay, là thành phải tránh xe hay xe phải tránh thành?''\\
\textit{(The boy replied: ``Since ancient times, should a fortress avoid the cart, or should the cart avoid the fortress?'')}

Khổng Tử ngẩn người rồi cúi đầu: ``Ta sai rồi. Mệnh lệnh xe đi vòng quanh thành!''\\
\textit{(Confucius was stunned, then bowed: ``I was wrong. Order the cart to go around the fortress!'')}

Học trò hỏi: ``Bậc thánh nhân sao lại nhường một đứa trẻ?''\\
\textit{(A disciple asked: ``Why does a sage yield to a child?'')}

Khổng Tử đáp: ``Chân lý không phân biệt tuổi tác. Ta sai là ta sai.''\\
\textit{(Confucius replied: ``Truth does not distinguish age. If I was wrong, I was wrong.'')}

\end{truyen}
\begin{baihoc}
  Khiêm tốn không phải yếu đuối; đó là sức mạnh để nhận ra và sửa sai bất kể chân lý đến từ đâu.\\
  \textit{(Humility is not weakness; it is the strength to recognize and correct mistakes regardless of where the truth comes from.)}
\end{baihoc}

\section{Socrates Và Lời Sấm Delphi}
\begin{truyen}{Socrates -- Người Khôn Nhất Vì Biết Mình Dốt}{Hy Lạp Cổ Đại}
\chuhoa{T}{}T hần Delphi phán truyền rằng Socrates là người thông minh nhất thành Athens.\\
\textit{(The Oracle of Delphi declared that Socrates was the wisest man in Athens.)}

Socrates bối rối vì tự biết mình không biết nhiều, liền đi gặp những người nổi tiếng học rộng.\\
\textit{(Socrates was puzzled, knowing his own ignorance, so he went to meet people renowned for their learning.)}

Ông phát hiện rằng các chính khách, nhà thơ, thợ thủ công đều tưởng mình biết những điều họ không thực sự biết.\\
\textit{(He discovered that politicians, poets, and craftsmen all believed they knew things they didn't actually know.)}

Socrates kết luận: ``Có lẽ ta thực sự khôn hơn họ một điều: ta biết rằng ta không biết.''\\
\textit{(Socrates concluded: ``Perhaps I am wiser than they are in one small way: I know that I do not know.'')}

Cả cuộc đời ông sống theo nguyên lý đó -- không ngừng đặt câu hỏi, không ngừng học hỏi.\\
\textit{(His whole life he lived by that principle -- never stopping questioning, never stopping learning.)}

\end{truyen}
\begin{baihoc}
  Biết mình không biết là cánh cửa đầu tiên của trí tuệ thực sự. Kẻ tưởng mình đã biết thì không còn chỗ để học thêm.\\
  \textit{(Knowing that you do not know is the first door of true wisdom. He who thinks he already knows has no room to learn more.)}
\end{baihoc}

\section{Warren Buffett Và Vòng Tròn Năng Lực}
\begin{truyen}{Vòng Tròn Năng Lực}{Warren Buffett -- Hiện Đại}
\chuhoa{W}{}W arren Buffett -- nhà đầu tư vĩ đại nhất thế giới -- luôn vẽ trên giấy một vòng tròn nhỏ.\\
\textit{(Warren Buffett, the world's greatest investor, always drew a small circle on paper.)}

Ông gọi đó là ``Vòng tròn năng lực'' -- phạm vi những gì ông thực sự hiểu rõ.\\
\textit{(He called it his ``Circle of Competence'' -- the domain of what he truly understood.)}

Khi ai hỏi về cổ phiếu công nghệ, ông thường đáp: ``Tôi không hiểu nó, tôi không đầu tư.''\\
\textit{(When asked about technology stocks, he usually replied: ``I don't understand it, so I don't invest in it.'')}

Trong khi nhiều người thua lỗ vì tham lam, Buffett làm giàu bằng cách biết rõ giới hạn của mình.\\
\textit{(While many lost money through greed, Buffett built wealth by clearly knowing his own limits.)}

Ông nói: ``Điều quan trọng không phải là vòng tròn to hay nhỏ, mà là biết chính xác ranh giới của nó.''\\
\textit{(He said: ``What matters is not how big or small the circle is, but knowing precisely where its boundaries are.'')}

\end{truyen}
\begin{baihoc}
  Sức mạnh thực sự nằm ở chỗ biết rõ ta giỏi gì và không giỏi gì -- và dũng cảm nói ``Tôi không biết.''\\
  \textit{(True strength lies in clearly knowing what you are good at and what you are not -- and the courage to say ``I don't know.'')}
\end{baihoc}

\section{Thiền Sư Và Chén Trà Đầy}
\begin{truyen}{Chén Trà Đã Đầy}{Nhật Bản -- Truyền Thuyết Thiền}
\chuhoa{M}{}M ột giáo sư đại học nổi tiếng đến gặp thiền sư Nan-in để học đạo.\\
\textit{(A famous university professor came to see Zen master Nan-in to learn the Tao.)}

Vị giáo sư liên tục nói về những gì ông đã biết, những học thuyết ông đã nghiên cứu.\\
\textit{(The professor kept talking about what he already knew, the theories he had already studied.)}

Nan-in lặng lẽ rót trà vào chén của khách và không dừng lại dù trà đã tràn ra ngoài.\\
\textit{(Nan-in quietly poured tea into the guest's cup and did not stop even as the tea overflowed.)}

Giáo sư kêu lên: ``Trà đầy rồi! Không thêm được nữa!''\\
\textit{(The professor cried out: ``The cup is full! No more can go in!'')}

Nan-in đáp: ``Ông cũng vậy. Ông đã đầy quan điểm và suy nghĩ. Làm sao tôi dạy được nếu ông không trống rỗng trước?''\\
\textit{(Nan-in replied: ``Like this cup, you are full of opinions and speculations. How can I teach you unless you first empty your cup?'')}

\end{truyen}
\begin{baihoc}
  Để tiếp nhận trí tuệ mới, trước hết phải làm trống chiếc chén kiến thức của mình. Học trò tốt nhất là người đến với tâm thế trống rỗng.\\
  \textit{(To receive new wisdom, one must first empty the cup of knowledge. The best student is one who comes with an empty mind.)}
\end{baihoc}

\section{Newton Và Bờ Biển Kiến Thức}
\begin{truyen}{Newton -- Đứa Trẻ Nhặt Sỏi Bên Bờ Đại Dương}{Nước Anh -- Thế Kỷ XVII}
\chuhoa{I}{}I saac Newton -- người khám phá ra lực hấp dẫn, người đặt nền móng vật lý cổ điển -- từng nói:

\textit{(Isaac Newton -- who discovered gravity and laid the foundation of classical physics -- once said:)}

``Tôi không biết mình trông như thế nào trước mắt thế giới, nhưng với chính mình, tôi chỉ như một đứa trẻ đang chơi trên bờ biển.''\\
\textit{(``I do not know what I may appear to the world, but to myself I seem to have been only like a boy playing on the seashore.'')}

``Đôi khi vui vẻ nhặt một hòn cuội nhẵn hơn bình thường, hoặc một vỏ sò đẹp hơn bình thường.''\\
\textit{(``Now and then diverting myself in finding a smoother pebble or a prettier shell than ordinary.'')}

``Trong khi đại dương vĩ đại của chân lý vẫn còn chưa được khám phá nằm trước tôi.''\\
\textit{(``While the great ocean of truth lay all undiscovered before me.'')}

Người vĩ đại nhất trong lịch sử khoa học vẫn xem mình chỉ là một đứa trẻ tò mò bên bờ đại dương.\\
\textit{(The greatest man in the history of science still saw himself as merely a curious child beside the ocean.)}

\end{truyen}
\begin{baihoc}
  Người học càng nhiều càng thấy mình biết ít. Khiêm tốn là dấu hiệu của trí tuệ thực sự.\\
  \textit{(The more one learns, the more one sees how little one knows. Humility is the hallmark of true wisdom.)}
\end{baihoc}
